\chapter{Mapa Lógico de Datos}

El mapa lógico resume cómo los datos transitan desde archivos raw heterogéneos hasta variables finales listas para validación, persistencia y análisis. En este proyecto, la secuencia observable es: contrato de lectura en metadata, lectura real por fuente, staging por fuente, integración sobre dimensión maestra, métricas derivadas y persistencia final en CSV y SQLite.

\begin{table}[htbp]
\centering
\footnotesize
\caption{Mapa lógico resumido desde origen hasta variable final.}
\begin{tabular}{|p{1.2cm}|p{4.0cm}|p{5.5cm}|p{3.4cm}|}
\hline
Fase/fuente & Origen & Transformación & Destino final \\
\hline
Todas & codigo\_comuna / Código / Código comuna & Parseo numérico y formateo string de 5 dígitos & codigo\_comuna \\
A & mmpqc\_2024 + mmpzc\_2024 & No Aplica -> 0; No Recepcionado -> NA; suma con min\_count=2 & areas\_verdes\_m2 \\
B & iadm41\_2024 & Renombre y conservación en miles de pesos nominales 2024 & ipp\_miles\_pesos \\
C & Porcentaje de personas en situación de pobreza por ingresos 2022 & Conversión de proporción 0-1 a porcentaje 0-100 y redondeo & pobreza\_ingresos\_pct \\
D & Población censada & Selección, tipificación Int64 y filtrado por universo & poblacion \\
Fase 6 & Constantes de referencia & Asignación explícita de años de referencia por variable & anio\_poblacion / anio\_pobreza / anio\_areas\_verdes / anio\_ingresos \\
Fase 6 & areas\_verdes\_m2 + poblacion & División controlada con población > 0 & areas\_verdes\_m2\_hab \\
Fase 6 & ipp\_miles\_pesos + poblacion & Conversión a pesos por habitante y control de infinitos & ipp\_pesos\_hab \\
\hline
\end{tabular}
\end{table}


El flujo anterior muestra dos criterios centrales. Primero, las variables finales solo se construyen a partir de campos realmente presentes en las fuentes y en los artefactos staging. Segundo, cualquier transformación semántica relevante queda documentada en el código y en los reportes de salida, por ejemplo la conversión de pobreza desde proporción a porcentaje o la suma de componentes de áreas verdes para producir \texttt{areas\_verdes\_m2}.

El detalle extendido de artefactos y destinos se deja en anexos para facilitar trazabilidad técnica sin sobrecargar el cuerpo principal del informe.
